\section{Referencial Teórico}

Diante do gargalo tecnológico provocado pela fragmentação de bases de dados geográficas e ambientais, as agências de controle carecem de ferramentas capazes de prover uma visão unificada e contínua do risco sistêmico, um estado mental de percepção, compreensão e projeção frequentemente denominado pela literatura como Consciência Situacional \cite{endsley1995}. Desta forma, a principal vocação do SIRDA é fornecer inteligência acionável para órgãos de controle ambiental e de infraestrutura, tendo essa meta teoricamente embasada pelo modelo seminal de Mica R. Endsley.

Em seu trabalho \textit{``Toward a theory of situation awareness in dynamic systems''}, Endsley estipula que a Consciência Situacional não se resume à simples coleta de dados, mas divide-se obrigatoriamente em três grandes níveis cognitivos e sistêmicos: Percepção (Nível 1), Compreensão (Nível 2) e Projeção (Nível 3). Na arquitetura do SIRDA, a ingestão de anomalias térmicas via satélite (Nível 1) é cruzada geoespacialmente com infraestruturas críticas para determinar o contexto de risco (Nível 2), culminando em simulações topológicas que projetam o impacto iminente de incêndios sobre as redes de transmissão (Nível 3). É essa taxonomia que justifica plenamente o emprego de tecnologias de processamento de fluxos em tempo real (\textit{Streaming}) e de painéis interativos de resposta rápida que alimentam a camada analítica do projeto.

Para materializar essa estrutura complexa de Consciência Situacional sem ceder ao caos operacional de arquiteturas legadas, o SIRDA fundamenta-se nos princípios do DataOps. Diferentemente de abordagens tradicionais de rotinas ETL (Extração, Transformação e Carga), a plataforma assimila o DataOps conforme conceituado por Julian Ereth \cite{ereth_dataops_2018}. Em seu artigo \textit{``DataOps: towards a definition''}, Ereth formaliza que esse paradigma transcende a mera automação, caracterizando-se pela incorporação sistemática das práticas estruturadas da Engenharia de Software ao ciclo de vida da Engenharia de Dados. 

Para a gestão de catástrofes ambientais, essa abordagem teórica mostra-se indispensável. Quando o tempo de resposta é crítico, não há margem para erros estruturais nas tabelas ou quedas silenciosas de integração. A teoria de Ereth aplica-se ao nosso caso de uso por meio da imposição de tipagens estritas nos contratos de dados e na orquestração de alto nível dos processos. Isso evita que falhas em fontes primárias propaguem dados sujos ou geometrias inválidas para o banco de dados analítico final. A utilização de sistemas OLAP para a tomada de decisão é amplamente embasada na literatura \cite{chaudhuri_olap_1997}, sendo que a plataforma propõe a adoção deste conceito para o processamento de eventos escaláveis com o Apache Druid \cite{yang_druid_2014}. Tais práticas garantem que os gestores de crise tomem decisões pautadas em informações íntegras e validadas.
